% Part: computability
% Chapter: computability-theory
% Section: no-universal-function

\documentclass[../../../include/open-logic-section]{subfiles}

\begin{document}

\olfileid{cmp}{thy}{nou}
\olsection{কোনো সার্বজনীন গণনসাধ্য অপেক্ষক নেই}

আংশিক গণনসাধ্য অপেক্ষকদের জন্য সার্বজনীন এমন একটি আংশিক গণনসাধ্য
অপেক্ষক থাকলেও, সর্বত্র-সংজ্ঞায়িত গণনসাধ্য অপেক্ষকদের জন্য সার্বজনীন
কোনো সর্বত্র-সংজ্ঞায়িত গণনসাধ্য অপেক্ষক নেই।
% BN-SRC-144 TARGET-CORRECTION-BEGIN
\par\smallskip\noindent\textnormal{\small\textbf{উৎস-সংশোধন (BN-SRC-144):} হিমায়িত ইংরেজি পাঠের প্রথম বাক্যে আংশিক গণনসাধ্য অপেক্ষকটিকে partial computable functions-এর জন্য “total” বলা হয়েছে। আগের সার্বজনীন-অপেক্ষক উপপাদ্য ও এই অংশের পরের বাক্য অনুসারে এখানে “universal” অভিপ্রেত; বাংলা পাঠে সেই অর্থই রাখা হয়েছে।} % BN-SRC-144 TARGET-CORRECTION-END

\begin{thm}
\ollabel{thm:no-univ}
কোনো সার্বজনীন গণনসাধ্য অপেক্ষক নেই। অন্যভাবে বললে, এমন যেকোনো অপেক্ষক
$\fn{Un}'(k, x)$ গণনসাধ্য নয়, যার ধর্ম হলো: $f(x)$ সর্বত্র-সংজ্ঞায়িত
গণনসাধ্য অপেক্ষক হলে এমন একটি স্বাভাবিক সংখ্যা~$k$ আছে, যাতে প্রত্যেক~$x$-এর
জন্য $f(x) = \fn{Un}'(k,x)$।
\end{thm}

\begin{proof}
প্রমাণটি একটি সরল কর্ণীকরণ: $\fn{Un}'(k,x)$ সর্বত্র-সংজ্ঞায়িত ও গণনসাধ্য
হলে
\[
d(x) = \fn{Un}'(x, x) + 1
\]
অপেক্ষকটিও সর্বত্র-সংজ্ঞায়িত ও গণনসাধ্য হতো। কিন্তু সংজ্ঞা অনুসারে
$d(k)$ এবং $\fn{Un}'(k,k)$ সমান নয়। অতএব প্রত্যেক $k$-এর জন্য $d(x)$
ও~$\fn{Un}'(k, x)$-এর মান অন্তত একটি~$x$-এ আলাদা, যথা $x = k$-এ।
\end{proof}

\begin{explain}
উপরের \olref[uni]{thm:univ-comp} দেখায় যে এই কর্ণীকরণ যুক্তিটি এড়ানো
যায়, তবে শুধু সার্বজনীন অপেক্ষকটিকে আংশিক রাখার বিনিময়ে। অর্থাৎ,
$\fn{Un}$ সর্বত্র-সংজ্ঞায়িত গণনসাধ্য অপেক্ষকদের জন্য সার্বজনীন, শুধু সেটি
নিজে সর্বত্র-সংজ্ঞায়িত নয়। আংশিক ক্ষেত্রে কর্ণীকরণ যুক্তিটি কাজ করে না।
\end{explain}

\begin{prob}
  \olref{thm:no-univ}-এর প্রমাণের কর্ণীকরণ যুক্তিটি আংশিক ক্ষেত্রে কেন
  কাজ করে না তা বুঝতে $f(x) \simeq \fn{Un}(x,x)+1$ অপেক্ষকটি বিবেচনা
  করুন। এটি কি আংশিক গণনসাধ্য? তা হলে এর একটি সূচক~$e$ আছে, অর্থাৎ
  $f(x) \simeq \fn{Un}(e,x)$।~$f(e)$ সম্বন্ধে কী বলা যায়?
\end{prob}


\end{document}
